\documentclass[../../../include/open-logic-section]{subfiles}

\begin{document}

\olfileid{sth}{ord-arithmetic}{add}
\olsection{جمعِ ترتیبی}

فرض کنید می‌خواهیم $\alpha$ و $\beta$ را جمع کنیم. می‌توانیم به‌سادگی
یک {نسخه} از $\beta$ را بی‌درنگ پس از یک نسخه از $\alpha$ بگذاریم.
(باید \emph{نسخه} برداریم، زیرا از
\olref[ordinals][basic]{ordinalsaresubsets} می‌دانیم که یا
$\alpha \subseteq \beta$ یا $\beta \subseteq \alpha$.) اثر شهودی این
کار آن است که نخست دنباله‌ای $\alpha$تایی از گام‌ها را طی کنیم و
\emph{سپس} دنباله‌ای $\beta$تایی را. دنبالهٔ حاصل خوش‌ترتیب خواهد
بود؛ بنابراین، بنا بر
\olref[ordinals][ordtype]{thmOrdinalRepresentation}، با یک عدد
ترتیبی (یکتا) یکریخت است. می‌توان آن عدد ترتیبی را \emph{جمعِ}
$\alpha$ و $\beta$ دانست.

این همان اندیشهٔ شهودیِ پشت جمع ترتیبی است. برای تعریف دقیق آن، از
ایدهٔ گرفتن \emph{نسخه‌هایی} از مجموعه‌ها آغاز می‌کنیم. در اینجا
برچسب‌های اختیاریِ $0$ و $1$ را به کار می‌بریم تا معلوم بماند هر
شیء از کجا آمده است:

\begin{defn}\ollabel{defdissum}
\emph{جمعِ مجزای} $A$ و $B$ برابر است با
$A \disjointsum B = (A\times \{0\}) \cup (B \times \{1\})$.
\end{defn}

اکنون ترتیبی را روی زوج‌های اعداد ترتیبی تعریف می‌کنیم:

\begin{defn}
برای هر اعداد ترتیبیِ $\alpha_1, \alpha_2, \beta_1, \beta_2$، می‌گوییم:
\begin{align*}
	\tuple{\alpha_1, \alpha_2} \rlexless \tuple{\beta_1, \beta_2}\text{ اگر و تنها اگر }&
	\text{یا $\alpha_2 \in \beta_2$}\\
	& \text{یا هم $\alpha_2 = \beta_2$ و هم $\alpha_1 \in \beta_1$}
\end{align*}
\end{defn}

این یک ترتیب \emph{واژه‌نامه‌ایِ معکوس} است، زیرا نخست بر پایهٔ
مؤلفهٔ دوم و سپس بر پایهٔ مؤلفهٔ نخست مرتب می‌کنید. به یاد آورید که
می‌خواستیم $\alpha \ordplus \beta$ را نوع ترتیبِ نسخه‌ای از $\alpha$
که نسخه‌ای از $\beta$ در پی آن می‌آید تعریف کنیم. برای این منظور
می‌گوییم:

\begin{defn}\ollabel{defordplus}
برای هر دو عدد ترتیبیِ $\alpha$ و $\beta$، جمع آن‌ها برابر است با
$\alpha \ordplus \beta =
\ordtype{\alpha \disjointsum \beta, \rlexless}$.
\end{defn}
\noindent
توجه کنید که در اینجا اندکی از نمادگذاری سوءاستفاده کرده‌ایم؛ به‌دقت
باید به‌جای «$\rlexless$» می‌نوشتیم
«$\Setabs{\tuple{x,y}\in(\alpha\disjointsum\beta)^2}{x \rlexless y}$».
بااین‌حال، برای اختصار در ادامه نیز همین سوءاستفادهٔ نمادی را پی
می‌گیریم.

نتیجهٔ زیر، همراه با
\olref[ordinals][ordtype]{thmOrdinalRepresentation}، تأیید می‌کند
که تعریف ما خوش‌صورت است:

\begin{lem}\ollabel{ordsumlessiswo}
$\tuple{\alpha \disjointsum \beta, \rlexless}$ برای هر دو عدد ترتیبیِ
$\alpha$ و $\beta$ یک خوش‌ترتیب است.
\end{lem}

\begin{proof}
آشکار است که $\rlexless$ روی $\alpha \disjointsum \beta$ متصل است.
برای نشان‌دادن خوش‌بنیادیِ آن، زیرمجموعهٔ ناتهیِ
$X \subseteq \alpha \disjointsum \beta$ را ثابت کنید. بگذارید $Y$
زیرمجموعه‌ای از $X$ باشد که مؤلفهٔ دومِ اعضایش تا جای ممکن کوچک است؛
یعنی
$Y = \Setabs{\tuple{\gamma, i} \in X}{(\forall \tuple{\delta, j} \in X)i \leq j}$.
اکنون عضوی از $Y$ را برگزینید که کوچک‌ترین مؤلفهٔ نخست را دارد.
\end{proof}
\noindent
پس تعریفی خوب و صریح از جمع ترتیبی داریم. واقعیت زیر نیز غافلگیرکننده
نیست (به یاد آورید که بنا بر
\olref[z][infinity-again]{defnomega}، $1 = \{0\}$):
\begin{prop}
برای هر عدد ترتیبیِ $\alpha$،
$\alpha \ordplus 1 = \ordsucc{\alpha}$.
\end{prop}

\begin{proof}
یکریختیِ $f$ را از
$\ordsucc{\alpha} = \alpha \cup \{\alpha\}$ به
$\alpha\disjointsum1 = (\alpha \times \{0\})
\cup (\{0\} \times \{1\})$ در نظر بگیرید که برای
$\gamma \in \alpha$ با $f(\gamma) = \tuple{\gamma, 0}$ و برای
$\alpha$ با $f(\alpha) = \tuple{0, 1}$ داده می‌شود.
\end{proof}
\noindent
افزون بر این، به‌آسانی می‌توان نشان داد که جمع از چند شرط بازگشتی
پیروی می‌کند:

\begin{lem}\ollabel{ordadditionrecursion}
برای هر دو عدد ترتیبیِ $\alpha, \beta$ داریم:
\begin{align*}
	\alpha\ordplus 0 &= \alpha\\
	\alpha \ordplus  (\beta\ordplus 1) &= (\alpha \ordplus  \beta) \ordplus  1\\
	\alpha  \ordplus  \beta &= \supstrict_{\delta < \beta}(\alpha \ordplus  \delta) && \text{اگر $\beta$ عدد ترتیبیِ حدی باشد}
\end{align*}
\end{lem}

\begin{proof}
حالت‌ها را یک‌به‌یک بررسی می‌کنیم؛ نخست:
\begin{align*}
	\alpha \ordplus  0
	&= \ordtype{\alpha \times \{0\}, \rlexless}\\
	&= \alpha\\
	\alpha \ordplus (\beta \ordplus  1)
	%&= \ordtype{(\alpha\times \{0\}) \cup (\ordtype{\beta \ordplus  1}\times \{1\}), \rlexless} \\
	&= \ordtype{(\alpha\times \{0\}) \cup (\ordsucc{\beta}\times \{1\}), \rlexless} \\
	&= \ordtype{(\alpha\times \{0\}) \cup (\beta \times \{1\}), \rlexless} \ordplus 1\\
	&= (\alpha \ordplus  \beta) \ordplus  1
\end{align*}
اکنون فرض کنید $\beta \neq \emptyset$ حدی باشد. اگر
$\delta < \beta$، آن‌گاه $\delta\ordplus 1 < \beta$ نیز برقرار است؛
پس $\alpha \ordplus \delta$ یک پارهٔ آغازینِ سره از
$\alpha \ordplus \beta$ است. بنابراین $\alpha \ordplus \beta$ یک
کران بالای اکید برای
$X = \Setabs{\alpha \ordplus \delta}{\delta < \beta}$ است. افزون بر
این، اگر $\alpha \leq \gamma < \alpha \ordplus \beta$، آشکار است که
برای یک $\delta < \beta$ داریم
$\gamma = \alpha \ordplus \delta$. پس
$\alpha \ordplus \beta =
\supstrict_{\delta< \beta}(\alpha\ordplus \delta)$.
\end{proof}

اما واقعیت چشمگیری در کار است. برای تعریف جمع ترتیبی می‌توانستیم
\emph{به‌جای آن} صرفاً قضیهٔ بازگشت ترامتناهی را به کار بریم و
معادلات بازگشتی را دقیقاً به صورتِ
\olref{ordadditionrecursion} مقرر کنیم (البته با
«$\ordsucc{\beta}$» به‌جای «$\beta \ordplus 1$»).

پس برای تعریف عمل‌ها روی اعداد ترتیبی دو راه متفاوت داریم. می‌توانیم
آن‌ها را \emph{ترکیبی} تعریف کنیم: یک مجموعهٔ خوش‌ترتیب را به‌صراحت
بسازیم و نوع ترتیبش را در نظر بگیریم. یا می‌توانیم آن‌ها را
\emph{بازگشتی} تعریف کنیم، فقط با مقررکردن معادلات بازگشت. با اجرای
درست، نتیجه یکسان است. قضیهٔ بازگشت ترامتناهی تضمین می‌کند که این
معادلات بازگشت اعداد ترتیبیِ \emph{یکتایی} را تعیین می‌کنند.

حساب ترتیبی از بسیاری جهات درست مانند جمع اعداد طبیعی رفتار می‌کند.
برای نمونه، می‌توانیم نتیجهٔ زیر را اثبات کنیم:

\begin{lem}\ollabel{ordinaladditionisnice}
اگر $\alpha, \beta, \gamma$ اعداد ترتیبی باشند، آن‌گاه:
\begin{enumerate}
	\item\ollabel{ordaddition1} اگر $\beta < \gamma$، آن‌گاه
	$\alpha \ordplus \beta < \alpha \ordplus \gamma$؛
	\item\ollabel{ordaddition2} اگر
	$\alpha \ordplus \beta = \alpha\ordplus \gamma$، آن‌گاه
	$\beta = \gamma$؛
	\item\ollabel{ordaddition3}
	$\alpha \ordplus (\beta \ordplus \gamma) =
	(\alpha \ordplus \beta) \ordplus \gamma$؛ یعنی جمع شرکت‌پذیر است؛
	\item\ollabel{ordaddition4} اگر $\alpha \leq \beta$، آن‌گاه
	$\alpha \ordplus \gamma \leq \beta \ordplus \gamma$.
\end{enumerate}
\end{lem}

\begin{proof}
\olref{ordaddition3} را اثبات می‌کنیم و بقیه را به‌عنوان تمرین
واگذار می‌کنیم. اثبات با استقرای ترامتناهیِ ساده بر $\gamma$ و با
استفاده از \olref{ordadditionrecursion} است. وقتی $\gamma = 0$:
\[
(\alpha \ordplus  \beta) \ordplus  0 = \alpha \ordplus  \beta  = \alpha \ordplus  (\beta \ordplus  0)
\]
وقتی $\gamma = \delta\ordplus 1$، به‌عنوان فرض استقرا بپذیرید که
$(\alpha \ordplus \beta) \ordplus \delta =
\alpha \ordplus (\beta \ordplus \delta)$؛ اکنون با سه بار استفاده
از \olref{ordadditionrecursion}:
\begin{align*}
	(\alpha \ordplus  \beta) \ordplus  (\delta \ordplus  1) & = ((\alpha \ordplus  \beta) \ordplus  \delta)\ordplus 1\\
	& = (\alpha \ordplus  (\beta \ordplus  \delta)) \ordplus  1\\
	& = \alpha \ordplus  ((\beta \ordplus  \delta)\ordplus 1)\\
	& = \alpha \ordplus  (\beta \ordplus  (\delta\ordplus 1))
\end{align*}
وقتی $\gamma$ عدد ترتیبیِ حدی است، به‌عنوان فرض استقرا بپذیرید که
اگر $\delta \in \gamma$، آن‌گاه
$(\alpha \ordplus \beta) \ordplus \delta =
\alpha \ordplus (\beta \ordplus \delta)$؛ اکنون:
\begin{align*}
	(\alpha \ordplus  \beta) \ordplus  \gamma & = \supstrict_{\delta < \gamma}((\alpha \ordplus  \beta) \ordplus  \delta) \\
	&= \supstrict_{\delta < \gamma}(\alpha \ordplus  (\beta \ordplus  \delta))\\
	&= \alpha \ordplus  \supstrict_{\delta < \gamma}(\beta \ordplus  \delta)\\
	& = \alpha \ordplus  (\beta \ordplus  \gamma)
\end{align*}
\end{proof}

\begin{prob}
باقیِ \olref[sth][ord-arithmetic][add]{ordinaladditionisnice} را
اثبات کنید.
\end{prob}

از این جهات، جمع ترتیبی باید بسیار آشنا به نظر برسد. اما یک تفاوت
بنیادی میان جمع ترتیبی و جمع اعداد طبیعی وجود دارد.

\begin{prop}\ollabel{ordsumnotcommute}
جمع ترتیبی {جابجایی‌پذیر نیست}؛
$1 \ordplus \omega = \omega < \omega \ordplus 1$.
\end{prop}

\begin{proof}
توجه کنید که
$1 \ordplus \omega = \supstrict_{n < \omega} (1 \ordplus n)
= \omega \in \omega \cup \{\omega\} = \ordsucc{\omega} = \omega
\ordplus 1$.
\end{proof}
\noindent
هرچند این امر ممکن است در آغاز غافلگیرکننده باشد، \emph{نباید چنین
باشد}. از یک سو، وقتی $1 \ordplus \omega$ را در نظر می‌گیرید، به نوع
ترتیبی می‌اندیشید که از گذاشتن یک عضو اضافی \emph{پیش از} همهٔ اعداد
طبیعی به دست می‌آید. با استدلالی همانند هتل هیلبرت در
\olref[sfr][infinite][hilbert]{sec}، از نظر شهودی این عضو نخستینِ
اضافی نباید نوع ترتیب کلی را تغییر دهد. از سوی دیگر، وقتی
$\omega \ordplus 1$ را در نظر می‌گیرید، به نوع ترتیبی می‌اندیشید که
از گذاشتن یک عضو اضافی \emph{پس از} همهٔ اعداد طبیعی حاصل می‌شود.
و آن موجودی یکسره متفاوت است!

\end{document}
